\section{Related Work}
\label{sec:related}

\todo[color=green]{Are there related work in terms of methodology?}

\textbf{Text-SPARQL Datasets.} Several curated benchmarks have driven progress in KGQA. The QALD challenge series~\cite{usbeck2019qald,usbeck2021qald} provides hand-crafted questions over DBpedia and Wikidata, while LC-QuAD~\cite{trivedi2017lc} and LC-QuAD 2.0~\cite{dubey2019lcquad} use template-based generation with human paraphrasing to scale dataset size. WikiSP~\cite{wikisp2023} sources questions from search logs to better reflect real user information needs. Spinach~\cite{liu2024spinach} curates questions from Wikidata's ``Request a Query'' forum, capturing authentic complexity beyond synthetic benchmarks. All of these resources target open-domain knowledge graphs such as Wikidata, DBpedia, and Freebase; domain-specific graphs, which lack equivalent breadth or community efforts, have no comparable Text-SPARQL coverage. \todo[color=green]{It seems obvious that they cannot be transferred.}

\textbf{KGQA via NL-SPARQL Translation.} KGQA has seen significant advances with the emergence of Large Language Models (LLMs), which have achieved state-of-the-art performance compared to prior deep learning and classical approaches \cite{dabramo-etal-2025-investigating,sgpt2022generative}. LLM-based approaches include chain-of-thought prompting for SPARQL generation \cite{zahera2024generating}, fine-tuning on text-SPARQL pairs \cite{sgpt2022generative,neural2022machine}, and in-context learning with retrieved examples \cite{emonet2024llm}. Agentic approaches such as GRASP~\cite{walter2025grasp} and Spinach~\cite{liu2024spinach} employ LLMs with tool-use for interactive graph exploration.\todo[color=green]{you mentioned these are not fully accurate.} All of these approache require or at least benefit from Text-SPARQL examples. QueryBloom is thus, complementary to these systems, as it can automatically provide the example sets they require for arbitrary RDFS-based knowledge graphs.

\textbf{Synthetic Data Generation.} Template-based approaches such as LC-QuAD~\cite{trivedi2017lc} generate questions from fixed SPARQL templates, but require manual template design per knowledge graph. Labate et al.~\cite{labate2025synthetic} recently proposed using LLMs to replace crowdsourced paraphrasing of LC-QuAD's template-generated questions, showing that synthetic rephrasings can match human quality for fine-tuning; however, their method still depends on LC-QuAD's existing template infrastructure and is limited to Wikidata. Deviating from template-based generation, KGQAGen~\cite{zhang2025kgqagen}, employs LLMs to generate at most one Wikidata-specific Text-SPARQL pair for each seed entity drawn from Wikipedia vital articles with the aim of improving factual quality issues observed in earlier benchmarks. While similar to QueryBloom in generating the Text-SPARQL pairs grounded on knowledge embedded in the graph, their method is designed specifically for Wikidata and uses Wikidata-specific APIs. Moreover, it generates questions strictly bounded by the number of seed entities given. In contrast, QueryBloom is designed to be applicable on any RDFS-based knowledge graph given its SPARQL endpoint and comprehensively generates Text-SPARQL pairs covering the schema underlying the knowledge graph while optionally requiring a single seed for the whole knowledge graph. The two methods could be viewed as complementary because Wikidata does not follow the RDFS standard required by QueryBloom and offers various resources, including seed entities and APIs, that make KGQAGen possible. In contrast, QueryBloom is designed to operate on any RDFS-based knowledge graph and does not assume access to such specific resources.